\chapter{Rectal Prolapse}

\section*{Definition}
Rectal prolapse is the circumferential full-thickness protrusion of the rectal wall beyond the anal verge. It is classified as internal (intussusception within the rectum), mucosal (only rectal mucosa protrudes), or complete/full-thickness (all layers of the rectal wall protrude through the anus).

\section*{What Happens in Your Body}
Weakness of the pelvic floor musculature, laxity of the rectal suspensory ligaments, and increased intra-abdominal pressure allow the rectum to intussuscept and protrude through the anal canal. Chronic straining causes stretching of the pudendal nerves and further muscle weakness.

\section*{Causes}
\begin{itemize}
  \item \textbf{Chronic constipation:} Prolonged straining during defecation (most common cause)
  \item \textbf{Age-related:} Pelvic floor weakness, laxity of the uterosacral ligaments (post-menopausal women most affected)
  \item \textbf{Parity:} Vaginal childbirth, especially multiple deliveries or prolonged second stage
  \item \textbf{Neurologic:} Pudendal neuropathy, cauda equina syndrome, spinal cord injury, multiple sclerosis
  \item \textbf{Other:} Cystic fibrosis (pediatric prolapse), parasitic infection (schistosomiasis, trichuriasis), pelvic tumors, prior anorectal surgery, malnutrition, chronic cough, Ehlers-Danlos syndrome
\end{itemize}

\section*{When to See a Doctor}
See your doctor if:
\begin{itemize}
  \item Tissue protrudes from the anus with defecation or Valsalva, requiring manual reduction
  \item Bleeding, mucus discharge, or fecal incontinence associated with prolapse
  \item Acute prolapse with incarceration (irreducible, edematous, painful)
  \item Strangulation: dark, dusky tissue, severe pain, or necrosis of the prolapsed segment
  \item New-onset prolapse in a child (evaluate for cystic fibrosis or parasitic infection)
\end{itemize}

\section*{Related Symptoms}
\begin{itemize}
  \item Sensation of a mass or bulge at the anus, fecal incontinence, mucus discharge, rectal bleeding, incomplete evacuation, constipation
\end{itemize}
\textbf{Important:} Rectal prolapse should be distinguished from mucosal prolapse and hemorrhoidal disease; full-thickness prolapse shows concentric mucosal rings, while hemorrhoids show radial groove pattern.

\section*{Self-Care / Home Management}
\begin{itemize}
  \item Manual reduction: lubricate prolapse, gentle steady pressure to return tissue through the anal canal
  \item Avoid straining during defecation; use stool softeners (docusate, psyllium) to ensure soft, regular bowel movements
  \item Avoid heavy lifting and prolonged standing
  \item Pelvic floor muscle exercises (Kegel exercises) to strengthen the levator ani
  \item Reduce prolapse-prone positions: squatting or using a footstool to mimic a squat during defecation
\end{itemize}

\section*{How Your Doctor Will Evaluate You}
\begin{itemize}
  \item Physical exam: inspect for protruding tissue in lithotomy, squatting, or seated on a commode
  \item Digital rectal exam to assess sphincter tone and identify other rectal pathology
  \item Proctoscopy or sigmoidoscopy to evaluate for associated pathology (solitary rectal ulcer, polyp, malignancy)
  \item Defecography (barium or MRI) to distinguish internal from complete prolapse and identify other pelvic floor disorders
  \item Anal manometry if fecal incontinence is present to assess sphincter function
  \item Cystic fibrosis sweat test or stool ova/parasite exam in children with new prolapse
\end{itemize}

\section*{Treatment Options}
\begin{itemize}
  \item Conservative therapy: dietary modification, stool softeners, pelvic floor physical therapy for mild or partial prolapse
  \item Surgical repair for symptomatic full-thickness prolapse: perineal approach (Altemeier perineal proctosigmoidectomy, Delorme mucosectomy) for older or high-risk people
  \item Abdominal approach (laparoscopic ventral rectopexy, sacral colpopexy) for younger people with better tissue quality
  \item Thiersch wire or anal encirclement procedure rarely used; reserved for the frailest people
  \item Management of associated pelvic floor disorders: cystocele, enterocele, rectocele often require concurrent repair
\end{itemize}

\section*{References}
\begin{thebibliography}{99}
\bibitem{rectal_prolapse:ref1} Tou S, Brown SR, Nelson RL. ``Surgery for complete rectal prolapse in adults.'' \textit{Cochrane Database of Systematic Reviews}, 2015;11:CD001758.
\bibitem{rectal_prolapse:ref2} Bordcianou L, Hicks CW, Kaiser AM, et al. ``Rectal prolapse: An overview of clinical features, diagnosis, and patient-specific management strategies.'' \textit{Journal of Gastrointestinal Surgery}, 2014;18(5):1059--1069.
\bibitem{rectal_prolapse:ref3} Madoff RD, Mellgren A. ``One hundred years of rectal prolapse surgery.'' \textit{Diseases of the Colon \& Rectum}, 1999;42(4):441--450.
\bibitem{rectal_prolapse:ref4} Varma M, Rafferty J, Buie WD. ``Practice parameters for the management of rectal prolapse.'' \textit{Diseases of the Colon \& Rectum}, 2011;54(11):1339--1346.
\bibitem{rectal_prolapse:ref5} Corman ML. \textit{Colon and Rectal Surgery}, 6th ed. Lippincott Williams \& Wilkins, 2013.
\bibitem{rectal_prolapse:ref6} Steele SR, Hull TL, Read TE, et al. \textit{The ASCRS Textbook of Colon and Rectal Surgery}, 3rd ed. Springer, 2016.
\end{thebibliography}
